\documentclass[11pt]{article}
\usepackage[T1]{fontenc}
\usepackage{lmodern,amsmath,amssymb,amsthm,mathtools,microtype}
\usepackage[a4paper,margin=27mm]{geometry}
\title{The complete \(\eta^{-2/3}\) strong-drive energy coefficient}
\author{Maciej Janowicz}
\date{}
\begin{document}
\maketitle

\section{Symbols and exact quadratic reduction}

The symbolic calculation uses
\[
n\in\mathbb N_0,\qquad \delta\in\mathbb R,\qquad
\lambda=r^{-1}>0,\qquad \epsilon=\eta^{-1/3}>0.
\]
In particular, no sign assumption is imposed on the detuning \(\delta\).
After \(a=b-r\), \(r^3+\delta r=\eta\), direct normal ordering gives
\[
\frac HV=-\frac32r^4-\delta r^2+(2r^2+\delta)b^\dagger b
+\frac{r^2}{2}(b^{\dagger2}+b^2)-r\mathcal O_3+\mathcal O_4,
\]
where
\[
\mathcal O_3=b^{\dagger2}b+b^\dagger b^2,\qquad
\mathcal O_4=\frac12b^{\dagger2}b^2.
\]
For \(A=2r^2+\delta\) and
\(\Omega=\sqrt{A^2-r^4}\), take \(b=C_\lambda c+S_\lambda c^\dagger\),
\[
C_\lambda^2=\frac12\left(\frac A\Omega+1\right),\quad
S_\lambda^2=\frac12\left(\frac A\Omega-1\right),\quad
2C_\lambda S_\lambda=-\frac{r^2}{\Omega}.
\]
The chosen branch has \(C_\lambda>0,S_\lambda<0\).  The anomalous
coefficient is
\[
AC_\lambda S_\lambda+\frac{r^2}{2}(C_\lambda^2+S_\lambda^2)=0,
\]
which directly audits the sign.  If
\[
C_0=\sqrt{\frac12+\frac1{\sqrt3}},\qquad
S_0=-\sqrt{-\frac12+\frac1{\sqrt3}},
\]
then
\begin{align*}
C_\lambda&=C_0\left[1+\delta\left(-\frac13+\frac{\sqrt3}{6}\right)
\lambda^2+\delta^2\left(\frac{17}{72}-\frac{\sqrt3}{9}\right)
\lambda^4+O(\lambda^6)\right],\\
S_\lambda&=S_0\left[1+\delta\left(-\frac13-\frac{\sqrt3}{6}\right)
\lambda^2+\delta^2\left(\frac{17}{72}+\frac{\sqrt3}{9}\right)
\lambda^4+O(\lambda^6)\right].
\end{align*}

\section{Recursive perturbation theory}

Define
\[
\mathcal K(\lambda)=\lambda^2(H/V-E_{\rm vac}^{(2)})
=\omega(\lambda)N-\lambda\mathcal O_3(\lambda)
+\lambda^2\mathcal O_4(\lambda),
\]
\[
\omega(\lambda)=\sqrt{3+4\delta\lambda^2+\delta^2\lambda^4}.
\]
Expanding the complete operators gives
\(\mathcal K=\sum_{j=0}^4\lambda^jK_j+O(\lambda^5)\).  Thus \(K_3\)
contains the \(\lambda^2\) correction to every factor in \(\mathcal O_3\),
and \(K_4\) contains the corresponding correction to every factor in
\(\mathcal O_4\), as well as
\(-\delta^2N/(6\sqrt3)\).

For intermediate normalization, put \(\psi^{(0)}=|n\rangle\).  Coefficient
comparison in the eigenvalue equation gives, recursively,
\[
\kappa_q=\left\langle n\left|
\sum_{j=1}^qK_j\psi^{(q-j)}\right.\right\rangle,
\]
\[
\langle m|\psi^{(q)}\rangle=
\frac{\left\langle m\left|
\sum_{j=1}^q(K_j-\kappa_j)\psi^{(q-j)}
\right.\right\rangle}{\sqrt3(n-m)},\qquad m\ne n.
\]
Algebraic ladder action retains all reachable states and annihilates
forbidden negative occupations.  It gives
\[
\kappa_1=\kappa_3=0,
\]
\[
\kappa_2=\frac{2\sqrt3}{3}\delta n
-\frac{n^2+n}{12}-\frac1{72},
\]
and
\begin{align*}
\kappa_4={}&-\frac{\sqrt3}{18}\delta^2n
+\frac{\delta n^2}{18}-\frac{\sqrt3\delta n}{9}
+\frac{\delta n}{18}-\frac{\sqrt3\delta}{18}+\frac{13\delta}{108}\\
&+\frac{5\sqrt3 n^3}{1296}
+\left(\frac{5\sqrt3}{864}-\frac1{18}\right)n^2
+\left(\frac{139\sqrt3}{2592}-\frac1{18}\right)n
-\frac5{108}+\frac{67\sqrt3}{2592}.
\end{align*}

\section{Independent finite-sum audit}

The second calculation uses the generalized fourth-order finite formula
directly.  With \(D_k=\sqrt3(n-k)\) and
\((K_j)_{ab}=\langle a|K_j|b\rangle\),
\begin{align*}
\kappa_4={}&(K_4)_{nn}
+\sum_{k\ne n}\frac{(K_1)_{nk}(K_3)_{kn}
+(K_3)_{nk}(K_1)_{kn}+(K_2)_{nk}(K_2)_{kn}}{D_k}\\
&+\sum_{k,l\ne n}\frac{
(K_1)_{nk}(K_1)_{kl}(K_2)_{ln}
+(K_1)_{nk}(K_2)_{kl}(K_1)_{ln}
+(K_2)_{nk}(K_1)_{kl}(K_1)_{ln}}{D_kD_l}\\
&+\sum_{k,l,m\ne n}\frac{
(K_1)_{nk}(K_1)_{kl}(K_1)_{lm}(K_1)_{mn}}{D_kD_lD_m}
-\kappa_2\sum_{k\ne n}\frac{(K_1)_{nk}(K_1)_{kn}}{D_k^2}.
\end{align*}
This includes finite-\(r\) vertex and frequency corrections, \(V_4^2\),
all three \(V_3^2V_4\) placements, \(V_3^4\), and the normalization
subtraction.  The allowed shifts make every sum finite.  Exact simplification
reproduces the recursive \(\kappa_2\) and \(\kappa_4\) above identically.

\section{Conversion to the drive variable}

With a separate \(\epsilon=\eta^{-1/3}\),
\[
r=\epsilon^{-1}-\frac{\delta}{3}\epsilon
+\frac{\delta^3}{81}\epsilon^5+O(\epsilon^7),
\]
which satisfies \(r^3+\delta r=\epsilon^{-3}+O(\epsilon^5)\).
Reconstructing
\[
\frac{E_n}{V}=E_{\rm vac}^{(2)}+r^2
\left(\sqrt3n+\lambda^2\kappa_2+\lambda^4\kappa_4\right)
\]
before this substitution yields
\begin{equation}
\boxed{\begin{aligned}
\frac{E_n}{V}={}&-\frac32\eta^{4/3}
+\left[\delta+\sqrt3(n+\tfrac12)-1\right]\eta^{2/3}\\
&+\frac{\delta(1-2\delta)}6-\frac{6n^2+6n+1}{72}
+K_n(\delta)\eta^{-2/3}+O(\eta^{-4/3}),
\end{aligned}}
\end{equation}
where
\begin{align}
K_n(\delta)=\frac1{2592}\bigl\{&
96\delta^3+(144\sqrt3n-288+72\sqrt3)\delta^2\nonumber\\
&+(144n^2-288\sqrt3n+144n-144\sqrt3+312)\delta\nonumber\\
&+10\sqrt3n^3+(15\sqrt3-144)n^2
+(139\sqrt3-144)n+67\sqrt3-120\bigr\}.
\label{eq:K}
\end{align}
The exact quadratic part contributes
\[
\frac{\delta^3}{27}+\frac{\sqrt3\delta^2n}{18}
-\frac{\delta^2}{9}+\frac{\sqrt3\delta^2}{36};
\]
the remainder of \eqref{eq:K} is the complete interaction contribution.

The spacing correction is
\[
K_{n+1}-K_n=\frac{
72\sqrt3\delta^2+(144n-144\sqrt3+144)\delta
+15\sqrt3n^2+(30\sqrt3-144)n+82\sqrt3-144}{1296},
\]
and the second-difference correction is
\[
K_{n+2}-2K_{n+1}+K_n
=\frac{48\delta+10\sqrt3n+15\sqrt3-48}{432}.
\]

\section{Status}
\begin{center}
\begin{tabular}{p{0.22\linewidth}p{0.68\linewidth}}
\hline
Exact & Displacement identity, Bogoliubov transformation, sign cancellation,
ladder actions, and equality of the two finite calculations.\\
Formal & Fixed-\(n\) Rayleigh--Schr\"odinger expansion through
\(\eta^{-2/3}\).\\
Conjectural & Its identification with zeros of the global Olver
compensated-to-essential coefficient.\\
Open & A rigorous operator remainder uniform in \(n\), the global canonical
domain chain, and the Olver connection-error estimate.\\ \hline
\end{tabular}
\end{center}
\end{document}
